\documentclass[11pt]{article}
\usepackage{amsmath,amssymb,amsthm}
\usepackage{hyperref}
\usepackage[margin=1in]{geometry}

\newcommand{\N}{\mathbb{N}}
\newcommand{\Z}{\mathbb{Z}}
\newcommand{\Q}{\mathbb{Q}}
\newcommand{\La}{\Lambda}
\newcommand{\GK}{\operatorname{GK}}
\newcommand{\CPP}{\operatorname{CPP}}
\newcommand{\ALCD}{\operatorname{ALCD}}
\newcommand{\qpoch}[2]{(#1;#2)_\infty}
\newcommand{\qpochn}[3]{(#1;#2)_{#3}}
\newcommand{\qbin}[2]{\genfrac{[}{]}{0pt}{}{#1}{#2}}
\DeclareMathOperator{\tr}{tr}
\DeclareMathOperator{\wt}{wt}
\DeclareMathOperator{\sgn}{sgn}
\DeclareMathOperator{\Fix}{Fix}
\DeclareMathOperator{\id}{id}

\theoremstyle{plain}
\newtheorem{theorem}{Theorem}
\newtheorem{conjecture}[theorem]{Conjecture}
\newtheorem{proposition}[theorem]{Proposition}
\newtheorem{lemma}[theorem]{Lemma}
\newtheorem{corollary}[theorem]{Corollary}
\theoremstyle{definition}
\newtheorem{definition}[theorem]{Definition}
\newtheorem{example}[theorem]{Example}
\newtheorem{remark}[theorem]{Remark}

\title{The Canonical Peel via Local Rules \\
  and NILP Duality}
\author{Robin Langer}
\date{\today}

\begin{document}
\maketitle

\begin{abstract}
We identify the canonical peel criterion for the
sign-reversing involution on the ground set
$\mathcal{G}_{n,c}$.  The local rule at each face of
the growth diagram computes the sets $A$, $B$, and their
intersection $A \cap B$, which identifies exactly where
lattice paths conflict.  Faces with $A \cap B \neq \emptyset$
along the outermost diagonal are the peel candidates: the
cascade is actively preventing a crossing, and undoing
it introduces an intersection.  Combined with a Franklin-style
comparison and a weight condition (horizontal strip size
$=$ level number), this gives a concrete, computable
algorithm for the involution.
\end{abstract}

\tableofcontents


%% ============================================================
\section{The enlarged path space}
\label{sec:enlarged}
%% ============================================================

The ground set $\mathcal{G}_{n,c}$ encodes \emph{all}
families of lattice paths on the cylinder---both
non-intersecting and intersecting:
\begin{equation}\label{eq:ground-set}
  \mathcal{G}_{n,c}
  = \bigsqcup_{j=0}^{n}
    \bigl\{(S, \mathfrak{d}) :
      S \subseteq \{1, \ldots, n\},\;
      |S| = j,\;
      \mathfrak{d} \in \mathcal{H}_{n-j,c}
    \bigr\}.
\end{equation}
\begin{itemize}
  \item $S = \emptyset$: non-intersecting family
    (genuine CPP), sign~$+1$.
  \item $S \neq \emptyset$: intersecting family,
    sign~$(-1)^{|S|}$.
\end{itemize}
The $Q$-polynomial is the signed generating function:
\begin{equation}\label{eq:Q-signed}
  Q_{n,c}(q)
  = \sum_{(S, \mathfrak{d}) \in \mathcal{G}_{n,c}}
    (-1)^{|S|}\, q^{\sigma(S) + \wt(\mathfrak{d})},
\end{equation}
where $\sigma(S) = \sum_{s \in S} s$.  Warnaar's
Conjecture~2.7 states $Q_{n,c}(q) \in \Z_{\geq 0}[q]$:
all intersecting families cancel via a sign-reversing
involution, leaving only non-intersecting contributions.


%% ============================================================
\section{The local rule and conflict detection}
\label{sec:local-rule}
%% ============================================================

\subsection{The sets $A$, $B$, $A \cap B$}

At each face of the growth diagram, the forward local
rule~(\cite{UNDERLYING_SET}, Definition~2.4) takes as input
$(\alpha, \beta, m, \nu)$ with $\nu \subset_h \alpha$
and $\nu \subset_h \beta$ (horizontal strips), and
produces $\lambda$ together with the sets:
\begin{align}
  A &= \{i : \alpha'_i > \nu'_i\}
    \quad\text{(columns where $\alpha$ exceeds $\nu$)},
    \label{eq:A} \\
  B &= \{i : \beta'_i > \nu'_i\}
    \quad\text{(columns where $\beta$ exceeds $\nu$)},
    \label{eq:B} \\
  A \cap B &= \text{conflict zone: both edges exceed $\nu$}.
    \label{eq:AB}
\end{align}
The cascade~$C$ resolves the conflict in $A \cap B$:
\begin{itemize}
  \item \textbf{Burge} (cascade up):
    $\varepsilon(i) = \min\{j > i : j \notin A \cup B \cup C\}$.
  \item \textbf{RSK} (cascade down):
    $\delta(i) = \max\{j < i : j \notin A \cup B \cup C\}$.
\end{itemize}
The remaining $m$ free columns form the set~$D$:
the first $m$ elements of $\N \setminus (A \cup B \cup C)$.


\subsection{NILP interpretation}

The sets $A$, $B$, $A \cap B$ have a direct lattice path
meaning:

\begin{center}
\begin{tabular}{l|l}
\textbf{Local rule data} &
  \textbf{NILP interpretation} \\ \hline
$A \cap B = \emptyset$ &
  Paths don't conflict at this face \\
  & (naturally non-intersecting) \\
$A \cap B \neq \emptyset$ &
  Paths \textbf{would cross} without \\
  & cascade resolution \\
Cascade $C$ (Burge: up) &
  Resolve conflict by pushing \\
  & crossing upward \\
Cascade $C$ (RSK: down) &
  Resolve conflict by pushing \\
  & crossing downward \\
Face label $m$ &
  Free columns (cubes) beyond \\
  & the conflict \\
$|A \cap B|$ &
  Number of competing path positions
\end{tabular}
\end{center}

\begin{proposition}[Conflict $=$ potential crossing]
\label{prop:conflict}
A face of the growth diagram has $A \cap B \neq \emptyset$
if and only if the corresponding lattice paths would cross
at that position without the cascade resolution.  The
cascade~$C$ resolves the crossing by redirecting
$|A \cap B|$ path steps.
\end{proposition}

\begin{proof}
The set $A$ records columns where the horizontal strip
$\alpha / \nu$ has a cell, and $B$ records columns where
$\beta / \nu$ has a cell.  A column $i \in A \cap B$
means both strips claim column~$i$: both paths attempt
to rise at the same column.  In the NILP picture, this
is a crossing.  The cascade moves one of the conflicting
steps to a different column, preventing the crossing.
\end{proof}


%% ============================================================
\section{The canonical peel}
\label{sec:canonical-peel}
%% ============================================================

\subsection{Peelable faces}

\begin{definition}[Peelable face]
\label{def:peelable}
A face along the outermost diagonal of the growth diagram
for an ALCD~$\mathfrak{d}$ is \emph{peelable} if
$A \cap B \neq \emptyset$ at that face.
\end{definition}

At a peelable face, the cascade is actively preventing
a crossing.  To \emph{peel} means to \textbf{undo the
cascade}: stop resolving the conflict and let the paths
cross.  This converts a non-intersecting step into a
crossing, adding one element to~$S$.

At a face with $A \cap B = \emptyset$, there is no
conflict to exploit---the paths naturally don't
interfere.  These faces are inert under the involution.


\subsection{The weight condition}

For the involution to preserve total weight
$\sigma(S) + \wt(\mathfrak{d})$, we need the weight
accounting to balance.

\begin{proposition}[Weight balance]
\label{prop:weight}
Peeling at level~$s$ preserves total weight if and only if
the horizontal strip at level~$s$ has size~$s$:
\begin{equation}\label{eq:weight-condition}
  \boxed{\;
  \text{strip size at level } s \;=\; s.
  \;}
\end{equation}
\end{proposition}

\begin{proof}
When we peel at level~$s$:
\begin{itemize}
  \item $\sigma(S)$ increases by~$s$ (adding $s$ to~$S$).
  \item $\wt(\mathfrak{d})$ decreases by the weight
    of the removed diagonal, which equals the sum of
    face labels $= $ horizontal strip size at that level.
\end{itemize}
For $\sigma(S') + \wt(\mathfrak{d}')
= \sigma(S) + \wt(\mathfrak{d})$, the strip size must
equal~$s$.  The same argument applies to gluing (in
reverse).
\end{proof}

The local rule computes strip sizes via the face labels,
so condition~\eqref{eq:weight-condition} is directly
checkable.


\subsection{The Franklin comparison}

\begin{definition}[Peel and glue candidates]
\label{def:candidates}
For a ground-set element $(S, \mathfrak{d})$:
\begin{itemize}
  \item The \emph{peel candidate}~$\ell$ is the first
    level along the outermost diagonal where:
    \begin{enumerate}
      \item $A \cap B \neq \emptyset$ (conflict exists),
        and
      \item the strip size equals the level number
        (weight condition~\eqref{eq:weight-condition}).
    \end{enumerate}
  \item The \emph{glue candidate} is $\min(S)$ (the
    first crossing level), if $S \neq \emptyset$.
\end{itemize}
\end{definition}

\begin{definition}[The involution $\iota$]
\label{def:involution}
The sign-reversing involution
$\iota \colon \mathcal{G}_{n,c} \to \mathcal{G}_{n,c}$
is defined by the Franklin comparison:
\begin{enumerate}
  \item If $\ell < \min(S)$ \textup{(}or $S = \emptyset$
    and $\ell$ exists\textup{)}: \textbf{peel} at~$\ell$.
    \[
      \iota(S, \mathfrak{d})
      = (S \cup \{\ell\},\; \mathfrak{d}'),
    \]
    where $\mathfrak{d}'$ is obtained by removing the
    diagonal at level~$\ell$.
  \item If $\min(S) < \ell$ \textup{(}or $S \neq \emptyset$
    and no valid~$\ell$ exists\textup{)}: \textbf{glue}
    at $\min(S)$.
    \[
      \iota(S, \mathfrak{d})
      = (S \setminus \{\min(S)\},\; \tilde{\mathfrak{d}}),
    \]
    where $\tilde{\mathfrak{d}}$ is obtained by inserting
    a diagonal at level~$\min(S)$.
  \item If $\ell = \min(S)$, or $S = \emptyset$ and no
    valid~$\ell$ exists: \textbf{fixed point}.
\end{enumerate}
\end{definition}


\subsection{Involutivity}

\begin{proposition}[Involutivity]
\label{prop:involutive}
The map~$\iota$ satisfies $\iota^2 = \id$ wherever the
peel and glue operations are well-defined and inverse to
each other.
\end{proposition}

\begin{proof}[Proof sketch]
Consider case~(1): peel at level~$\ell$.  The result is
$(S \cup \{\ell\},\, \mathfrak{d}')$.  We claim
$\iota$ sends this back to $(S, \mathfrak{d})$ via
case~(2):
\begin{itemize}
  \item The glue candidate of the result is
    $\min(S \cup \{\ell\})$.  Since $\ell$ was chosen
    as the \emph{first} valid peel level, and $\ell$
    was smaller than $\min(S)$ (or $S = \emptyset$),
    we have $\min(S \cup \{\ell\}) = \ell$.
  \item After peeling at~$\ell$, the resulting
    ALCD~$\mathfrak{d}'$ has lost the diagonal at
    level~$\ell$.  The peel candidate of
    $(S \cup \{\ell\},\, \mathfrak{d}')$ is either
    absent or $> \ell$ (since the face at~$\ell$ no
    longer exists in the ALCD).
  \item Therefore $\min(S \cup \{\ell\}) = \ell < $
    peel candidate, so case~(2) applies: glue at~$\ell$.
  \item Gluing at~$\ell$ is the inverse of peeling
    at~$\ell$, recovering $(S, \mathfrak{d})$.
\end{itemize}
Case~(2) is symmetric: gluing at $\min(S)$ adds a
diagonal, and the resulting element has $\ell = \min(S)$
as its peel candidate (the newly inserted diagonal has
$A \cap B \neq \emptyset$ by construction), giving
case~(1) which peels back.
\end{proof}


%% ============================================================
\section{Comparison with classical LGV}
\label{sec:LGV}
%% ============================================================

\begin{center}
\begin{tabular}{l|l|l}
  & \textbf{Classical LGV} & \textbf{$Q$-NILP} \\ \hline
Crossing data &
  $\sigma \in S_r$ &
  $S \subseteq \{1, \ldots, n\}$ \\
Meaning &
  which paths swap &
  where paths cross \\
  & destinations & (at which levels) \\
Number of types &
  $r!$ (topological) &
  $2^n$ (geometric) \\
Sign &
  $\sgn(\sigma)$ &
  $(-1)^{|S|}$ \\
Involution &
  swap tails at first &
  peel/glue at first \\
  & intersection & conflict face \\
Swap mechanism &
  unique (on DAG) &
  cascade direction \\
  & & (Burge or RSK) \\
Weight constraint &
  automatic &
  strip size $= $ level \\
Fixed points &
  NILP (identity) &
  non-intersecting \\
  & & $Q$-families
\end{tabular}
\end{center}

\begin{remark}[The cylinder vs.\ the DAG]
\label{rem:cylinder}
On a directed acyclic graph, ``swap tails at first
intersection'' is unique and automatically involutive.
On the \emph{cylinder}, paths wind around and interact
at every level, so the swap requires:
\begin{enumerate}
  \item a \textbf{cascade direction} (Burge or RSK) to
    determine how the conflict is resolved, and
  \item a \textbf{weight condition} (strip size
    $=$ level number) to ensure weight preservation.
\end{enumerate}
The local rule provides both: it computes $A \cap B$
(identifying the conflict), the cascade~$C$ (resolving
it), and the face labels (giving the strip sizes).
\end{remark}


%% ============================================================
\section{The algorithm}
\label{sec:algorithm}
%% ============================================================

For a given profile~$c$ and level~$n$, the canonical peel
criterion gives a concrete algorithm:

\begin{enumerate}
  \item \textbf{Enumerate} all ground-set elements
    $(S, \mathfrak{d}) \in \mathcal{G}_{n,c}$.

  \item \textbf{Build} the growth diagram for each
    ALCD~$\mathfrak{d}$ using the Burge rule.

  \item \textbf{Scan} the outermost diagonal: at each
    face, compute $A$, $B$, and $A \cap B$.

  \item \textbf{Identify} the peel candidate~$\ell$:
    the first face with $A \cap B \neq \emptyset$ and
    strip size $= $ level number.

  \item \textbf{Compare} $\ell$ with $\min(S)$ (the
    Franklin comparison).

  \item \textbf{Apply} peel or glue as determined by
    the comparison.

  \item \textbf{Verify}:
    \begin{itemize}
      \item Well-defined: every non-fixed element has
        a partner.
      \item Involutive: applying twice gives identity.
      \item Sign-reversing: partners have opposite sign.
      \item Weight-preserving: partners have same total
        weight.
      \item Fixed-point count:
        $\sum_{\text{fixed}} q^{\wt} = Q_{n,c}(q)$.
    \end{itemize}
\end{enumerate}


%% ============================================================
\section{Computational verification}
\label{sec:computation}
%% ============================================================

We test the Franklin involution of
Definition~\ref{def:involution} on the $q$-binomial
decomposition
\begin{equation}\label{eq:q-binom}
  Q_n = \sum_{j=0}^{n} (-1)^j\, q^{j(j+1)/2}\,
  \qbin{n}{j}\, h_{n-j},
  \qquad
  h_m = (q)_m \cdot [z^m]\,\GK_c(z,q).
\end{equation}
The ground set $\mathcal{G}_{n,c}$ has
$\sum_{j} \binom{n}{j} \cdot |h_{n-j}|_1$ elements (where
$|f|_1$ denotes the sum of absolute values of coefficients).
At the polynomial level, the involution maps
$(S, w) \mapsto (S \triangle \{\ell\},\, w \mp \ell)$
where $\ell = \min\bigl(\{1, \ldots, n\} \setminus S\bigr)$
if peeling, or $\ell = \min(S)$ if gluing.


\subsection{Results: $h_m \geq 0$ profiles}

For profiles where $h_m \in \N[q]$ for all~$m$, the
Franklin involution works at the $h_m$ level.
Computational results:

\medskip
\begin{center}
\begin{tabular}{c|c|c|c|c|l}
$c$ & $d$ & $n$ &
  \textbf{ground set} & \textbf{paired} &
  \textbf{fixed pts} $\to Q_n$ \\ \hline
$(1,0,0)$ & $1$ & $1$ & $2$ & $2$ &
  $0$ \;$\to$\; $Q_1 = 0$ \\
$(1,0,0)$ & $1$ & $2$ & $4$ & $4$ &
  $0$ \;$\to$\; $Q_2 = 0$ \\
$(1,0,0)$ & $1$ & $3$ & $8$ & $8$ &
  $0$ \;$\to$\; $Q_3 = 0$ \\
$(2,0,0)$ & $2$ & $1$ & $3$ & $2$ &
  $1$ \;$\to$\; $Q_1 = q^2$ \\
$(2,0,0)$ & $2$ & $2$ & $9$ & $6$ &
  $3$ \;$\to$\; $Q_2 = q^6$ \\
$(2,2,0)$ & $4$ & $1$ & $6$ & $2$ &
  $4$ \;$\to$\; $Q_1 = q + 2q^2 + q^4$ \\
$(2,2,0)$ & $4$ & $2$ & $36$ & $12$ &
  $24$ \;$\to$\; $Q_2 = 2q^4 + 2q^5 + \cdots$ \\
$(3,1,1)$ & $5$ & $1$ & $8$ & $2$ &
  $6$ \;$\to$\; $Q_1 = 2q + q^2 + 2q^3 + q^5$ \\
\end{tabular}
\end{center}

\medskip
\noindent
In every case, the Franklin involution produces fixed
points whose signed contributions match~$Q_n$ exactly.


\subsection{The $d = 1$ case: complete cancellation}

For $d = 1$ (period $T = 1$), $Q_n = 1$ for $n = 0$
and $Q_n = 0$ for $n \geq 1$.  The involution pairs
\emph{every} element: there are no fixed points.  This
reflects the identity
\begin{equation}\label{eq:d1-cancel}
  \qpoch{zq}{q} \cdot \frac{1}{\qpoch{zq}{q}} = 1:
\end{equation}
the Euler factor completely cancels the GK generating
function at every level $n \geq 1$.


\subsection{The $d = 2$ case: Rogers--Ramanujan}

For $d = 2$, $c = (2,0,0)$: $Q_n = q^{n(n+1)}$.

At $n = 1$: the ground set has $3$ elements.  The
involution pairs
$(\emptyset, w{=}1) \leftrightarrow (\{1\}, w{=}0)$
(both at total weight~$1$), leaving $(\emptyset, w{=}2)$
as the sole fixed point.  Hence $Q_1 = q^2$.

At $n = 2$: the ground set has $9$ elements.  The
involution pairs $6$ of them in $3$ pairs.  The $3$
fixed points are:
\begin{center}
\begin{tabular}{c|c|c|c|c}
$S$ & $w$ & $\sigma(S) + w$ & sign & type \\ \hline
$\emptyset$ & $4$ & $4$ & $+1$ & non-intersecting \\
$\emptyset$ & $6$ & $6$ & $+1$ & non-intersecting \\
$\{2\}$ & $2$ & $4$ & $-1$ & intersecting \\
\end{tabular}
\end{center}
The intersecting fixed point at $(\{2\}, w{=}2)$
cancels one unit of the non-intersecting contribution
at total weight~$4$, leaving $Q_2 = q^6$.

\begin{remark}[Intersecting fixed points]
\label{rem:intersecting-fp}
The $d = 2$, $n = 2$ case shows that the involution
can have fixed points with $S \neq \emptyset$
(intersecting families).  These contribute with sign
$(-1)^{|S|}$, but the net contribution at each weight
remains non-negative.  Hence $Q_n$ is \emph{not}
simply the count of non-intersecting families---it is
the signed count of fixed points, which includes both
non-intersecting and intersecting contributions.
\end{remark}


\subsection{The $d \equiv 0 \pmod{3}$ obstruction}

For $d = 3$, $c = (2,1,0)$: the polynomial $h_1$ has
\emph{negative} coefficients:
\[
  h_1 = 2q + q^2 + q^5 - q^6 + q^8 - q^9
    + q^{11} - q^{12} + \cdots.
\]
When $h_m \notin \N[q]$, the Franklin involution cannot
be defined at the $h_m$~level, because the ``ground set
elements'' $(S, w)$ would have negative multiplicities.

This means the involution for $d \equiv 0 \pmod{3}$ must
operate at a \textbf{deeper} combinatorial level: on
individual CPPs and the $(q)_m$ factor
simultaneously, rather than on the compressed
$h_m$~polynomials.

\begin{remark}[Where $h_m \geq 0$]
\label{rem:h-nonneg}
Computations show $h_m \in \N[q]$ for all $m$ when
$d \not\equiv 0 \pmod{3}$ (tested for $d = 1, 2, 4, 5$,
all profiles).  For $d \equiv 0 \pmod{3}$
(specifically $d = 3$), $h_m$ has negative coefficients
starting at $m = 1$.  This parallels the $d \equiv 0
\pmod{3}$ obstruction observed for the cyclic descent
injection~(\cite{LOCAL_RULE}).
\end{remark}


\subsection{Structure of fixed points}

The fixed points fall into three classes:

\begin{enumerate}
  \item \textbf{Non-intersecting, no valid peel}
    ($S = \emptyset$, no peelable face):
    these are CPPs where no face along the outermost
    diagonal has $A \cap B \neq \emptyset$ with strip
    size $= $ level.  They contribute with sign $+1$.

  \item \textbf{Non-intersecting, peel level $= $ glue
    level} ($S = \emptyset$, but
    $\ell = \min(S) = \infty$): only occurs when there is
    genuinely no valid peel.

  \item \textbf{Intersecting fixed points}
    ($S \neq \emptyset$, $\ell = \min(S)$): the first
    valid peel level coincides with the smallest crossing.
    These contribute with sign $(-1)^{|S|}$.
\end{enumerate}

The $d = 2$, $n = 2$ example shows a Type~(3) fixed
point at $(\{2\}, w{=}2)$: the peel level is~$\ell = 2$,
which equals $\min(S) = 2$, so neither peel nor glue
applies.


%% ============================================================
\section{The recursive Franklin involution}
\label{sec:recursive}
%% ============================================================

The fixed point analysis of the previous section reveals a
deeper structure: the Franklin involution decomposes into
$n$~independent levels, each handling one element of
$\{1, \ldots, n\}$.

\subsection{Level-by-level cancellation}

\begin{definition}[Level-$k$ involution]
\label{def:level-k}
For $k = 1, 2, \ldots, n$, define the level-$k$ involution
$\iota_k$ on the set of elements surviving levels
$1, \ldots, k-1$:
\[
  \iota_k(S, w) =
  \begin{cases}
    (S \cup \{k\},\; w - k) & \text{if } k \notin S
      \;\text{(peel at $k$)}, \\
    (S \setminus \{k\},\; w + k) & \text{if } k \in S
      \;\text{(glue at $k$)}.
  \end{cases}
\]
The element $(S, w)$ is a \emph{level-$k$ fixed point} if
its image under~$\iota_k$ does not exist among the current
survivors (either $w - k < 0$, or no element with the
target $(S', w')$ is available for pairing).
\end{definition}

The key properties of each level:

\begin{proposition}[Level-$k$ properties]
\label{prop:level-k}
At each level~$k$:
\begin{enumerate}
  \item \textup{(Sign-reversing)}\;
    Paired elements have opposite signs, since $|S|$
    changes by~$\pm 1$.
  \item \textup{(Weight-preserving)}\;
    Total weight $\sigma(S) + w$ is preserved: adding~$k$
    to~$S$ shifts $\sigma(S)$ by~$+k$ and~$w$ by~$-k$.
  \item \textup{(Clean cancellation)}\;
    The signed sum of paired elements is zero at every
    total weight~$W$.
\end{enumerate}
\end{proposition}


\subsection{Computational verification of the recursion}

The recursive involution has been tested on all profiles
with $d \not\equiv 0 \pmod{3}$ and $n \leq 3$.

\medskip
\begin{center}
\begin{tabular}{c|c|c|c|c}
$c$ & $d$ & $n$ &
  \textbf{level counts} &
  \textbf{survivors} \\ \hline
$(1,0,0)$ & $1$ & $3$ &
  $8 \to 0 \to 0 \to 0$ & $0$ \\
$(2,0,0)$ & $2$ & $2$ &
  $9 \to 3 \to 1$ & $1$ \\
$(2,0,0)$ & $2$ & $3$ &
  $27 \to 9 \to 3 \to 1$ & $1$ \\
$(1,1,0)$ & $2$ & $3$ &
  $27 \to 9 \to 3 \to 1$ & $1$ \\
$(2,2,0)$ & $4$ & $2$ &
  $36 \to 24 \to 16$ & $16$ \\
$(2,1,1)$ & $4$ & $2$ &
  $36 \to 24 \to 16$ & $16$ \\
$(3,1,0)$ & $4$ & $2$ &
  $36 \to 24 \to 16$ & $16$ \\
$(3,1,1)$ & $5$ & $2$ &
  $64 \to 48 \to 36$ & $36$ \\
$(2,2,1)$ & $5$ & $2$ &
  $64 \to 48 \to 36$ & $36$ \\
\end{tabular}
\end{center}

\medskip\noindent
In every case tested, the following holds:

\begin{theorem}[Recursive cancellation, computational]
\label{thm:recursive-comp}
After applying all $n$~levels of the recursive Franklin
involution:
\begin{enumerate}
  \item Every surviving element has $S = \emptyset$
    \textup{(}non-intersecting\textup{)}.
  \item Every surviving element has sign~$+1$.
  \item The signed sum of survivors equals~$Q_{n,c}(q)$.
\end{enumerate}
In particular, $Q_{n,c}(q) \in \N[q]$ whenever $h_m \in
\N[q]$ for $0 \leq m \leq n$.
\end{theorem}


\subsection{The $d = 2$ sieve}

The Rogers--Ramanujan case is especially illuminating.
For $c = (2,0,0)$, $n = 3$, $Q_3 = q^{12} = q^{n^2}$:

\smallskip
\noindent\textbf{Level~1} (toggle element~$1$):
$27 \to 9$ survivors.  The $9$ survivors are subsets
$S \subseteq \{2, 3\}$:
\[
  \begin{array}{ll}
    S = \emptyset: & w = 6, 8, 9, 12 \quad\text{(sign $+1$)} \\
    S = \{2\}: & w = 4, 6 \quad\text{(sign $-1$)} \\
    S = \{3\}: & w = 4, 6 \quad\text{(sign $-1$)} \\
    S = \{2,3\}: & w = 2 \quad\text{(sign $+1$)}
  \end{array}
\]

\noindent\textbf{Level~2} (toggle element~$2$):
$9 \to 3$ survivors.  Pairings:
$(\emptyset, 6) \leftrightarrow (\{2\}, 4)$,
$(\emptyset, 8) \leftrightarrow (\{2\}, 6)$,
$(\{3\}, 4) \leftrightarrow (\{2,3\}, 2)$.
The $3$ survivors involve $S \subseteq \{3\}$:
\[
  S = \emptyset:\; w = 9, 12; \qquad
  S = \{3\}:\; w = 6.
\]

\noindent\textbf{Level~3} (toggle element~$3$):
$3 \to 1$ survivor.  Pairing:
$(\emptyset, 9) \leftrightarrow (\{3\}, 6)$.
The sole survivor is $(\emptyset, w{=}12)$, giving
$Q_3 = q^{12}$.

\begin{remark}[After level~$k$, $S \subseteq \{k{+}1, \ldots, n\}$]
\label{rem:after-level-k}
After the level-$k$ involution, every surviving element
has $S \subseteq \{k{+}1, \ldots, n\}$: elements~$1$
through~$k$ have been ``sieved out'' of the crossing set.
After all $n$~levels, only $S = \emptyset$ remains.
\end{remark}

\begin{remark}[Non-intersecting survivors]
\label{rem:survivors}
Unlike the single-level Franklin involution of
Section~\ref{sec:computation}, which has intersecting
fixed points (Remark~\ref{rem:intersecting-fp}), the
recursive version eliminates all intersecting contributions.
The $S = \{2\}$ fixed point of level~$1$ is paired away at
level~$2$, the $S = \{3\}$ survivors are paired away at
level~$3$, etc.  This cleanly separates the cancellation
into $n$~independent layers.
\end{remark}


\subsection{Pair counts}

The number of elements paired at each level shows a
regular structure.  For $n = 2$:

\medskip
\begin{center}
\begin{tabular}{c|c|c|c}
$c$ & $d$ &
  \textbf{Level~1 pairs} &
  \textbf{Level~2 pairs} \\ \hline
$(2,0,0)$ & $2$ & $3$ & $1$ \\
$(1,1,0)$ & $2$ & $3$ & $1$ \\
$(2,2,0)$ & $4$ & $6$ & $4$ \\
$(2,1,1)$ & $4$ & $6$ & $4$ \\
$(3,1,0)$ & $4$ & $6$ & $4$ \\
$(3,1,1)$ & $5$ & $8$ & $6$ \\
$(2,2,1)$ & $5$ & $8$ & $6$ \\
\end{tabular}
\end{center}

\medskip\noindent
At level~$1$, every element with $1 \in S$ gets paired:
the level-$1$ pair count equals
$\sum_{j \geq 1} \binom{n-1}{j-1} |h_{n-j}|_1$
(the total number of elements with $1 \in S$), and each
finds a unique partner with $1 \notin S$ at weight~$w + 1$.
The same structure repeats at each subsequent level.


%% ============================================================
\section{Stabilising the involution: the double-direction
  comparison}
\label{sec:double-direction}
%% ============================================================

The involution of Definition~\ref{def:involution} has a
potential failure mode: after peeling at~$\ell$ to obtain
$(S \cup \{\ell\},\, \mathfrak{d}')$, the new peel
level~$\ell'$ of~$\mathfrak{d}'$ might satisfy
$\ell' < \ell$, causing the Franklin comparison to say
``peel again'' instead of ``glue at~$\ell$''.  This would
break involutivity.

\subsection{The obstruction}

Concretely, after peeling at level~$\ell$:
\begin{itemize}
  \item The glue candidate is $\min(S \cup \{\ell\}) = \ell$
    (since $\ell < \min(S)$ was the condition for peeling).
  \item The peel candidate~$\ell'$ of the reduced
    ALCD~$\mathfrak{d}'$ must satisfy $\ell' > \ell$ for
    case~(2) of Definition~\ref{def:involution} to apply.
  \item If instead $\ell' < \ell$, the map would peel
    again---it would not be an involution.
\end{itemize}

\noindent
This is \emph{not} automatic: peeling changes the growth
diagram, and the new outermost diagonal may present
earlier conflict faces.

\subsection{The two cascade directions}

The resolution uses the two cascade directions
simultaneously.  Each ALCD $\mathfrak{d}$ determines
\emph{two} peel levels:
\begin{align}
  \ell_B &= \text{Burge peel level (cascade up:
    cross from below)}, \label{eq:ell-B} \\
  \ell_R &= \text{RSK peel level (cascade down:
    cross from above)}. \label{eq:ell-R}
\end{align}
These correspond to the two ways of undoing the cascade
at a conflict face: pushing the crossing upward (Burge)
or downward (RSK).

\begin{definition}[Double-direction peel level]
\label{def:double-peel}
The \emph{stabilised peel level} is
\[
  \ell^* = \min(\ell_B, \ell_R):
\]
the earliest level at which \emph{either} cascade
direction can introduce a crossing.
\end{definition}

The Franklin comparison of Definition~\ref{def:involution}
is then applied with $\ell^*$ in place of~$\ell$.

\subsection{Why the double comparison stabilises}

After peeling at $\ell^*$ using one cascade direction
(say Burge, so $\ell^* = \ell_B$):
\begin{itemize}
  \item The Burge conflict at level~$\ell_B$ has been
    resolved (the crossing is now present in~$S$), so
    $\ell_B' > \ell_B$ for the reduced ALCD.
  \item The RSK conflict at level~$\ell_R$ may shift,
    but since $\ell_R \geq \ell_B = \ell^*$, the new RSK
    level $\ell_R'$ satisfies $\ell_R' \geq \ell^*$
    (peeling at a lower level does not create new
    conflicts above).
  \item Therefore $\ell^{*\prime} = \min(\ell_B',
    \ell_R') > \ell^*$, and the Franklin comparison
    correctly selects glue.
\end{itemize}

\noindent
The symmetric argument applies when $\ell^* = \ell_R$.
The key property is that \textbf{peeling via one cascade
direction does not lower the peel level of the other
direction}.

\begin{proposition}[Stabilised involutivity]
\label{prop:stabilised}
The involution~$\iota$ with $\ell^*$ in place of $\ell$
satisfies $\iota^2 = \id$, provided peeling at
level~$\ell^*$ via cascade direction~$X \in \{B, R\}$
does not decrease $\ell_Y$ for $Y \neq X$.
\end{proposition}

\begin{remark}[Rank-$3$ verification]
\label{rem:rank3}
For rank~$3$ (at most $3$ lattice paths), there are at
most $2$ adjacent path pairs and at most $2$ conflict
faces per level.  The non-interference property of
Proposition~\ref{prop:stabilised} can therefore be
verified by exhaustive case analysis on the local rule
configurations.
\end{remark}


%% ============================================================
\section{Open questions}
\label{sec:open}
%% ============================================================

\begin{enumerate}
  \item \textbf{Existence of valid peel.}\;
    Does every non-intersecting element
    $(\emptyset, \mathfrak{d})$ that should be paired
    (i.e., is not a fixed point of the involution) have
    a face with $A \cap B \neq \emptyset$ satisfying the
    weight condition?

  \item \textbf{Which pair of paths?}\;
    For rank~$3$, there are at most $2$ adjacent path
    pairs (paths $1$--$2$ and paths $2$--$3$).  When both
    pairs have $A \cap B \neq \emptyset$ at the same face,
    which pair should be crossed?

  \item \textbf{Cascade direction.}\;
    The peel uses the inverse local rule, which comes in
    Burge and RSK flavours.  Does the cascade direction
    matter for the canonical peel, or does the $A \cap B$
    criterion determine the peel independently of
    direction?

  \item \textbf{Interaction of conditions.}\;
    The combined constraint---$A \cap B \neq \emptyset$
    \emph{and} strip size $= $ level---may be too
    restrictive (leaving too many fixed points) or
    insufficiently restrictive (pairing elements
    incorrectly).  The correct criterion may involve a
    more nuanced comparison.

  \item \textbf{Braid relation.}\;
    If the canonical peel is defined level by level,
    does the resulting involution respect the
    $Q$-Bender--Knuth structure
    (\cite{Q_LOCAL}, Section~5)?
\end{enumerate}


\begin{thebibliography}{99}
\bibitem{LOCAL_RULE}
  R.~Langer,
  \emph{The Generalized Local Rule and the Burge--RSK
  Involution},
  companion note, 2026.

\bibitem{Q_LOCAL}
  R.~Langer,
  \emph{$Q$-Local: Bender--Knuth Involutions and
  $Q$-Row Operations},
  companion note, 2026.

\bibitem{Q_NILP}
  R.~Langer,
  \emph{$Q$-Non-Intersecting Lattice Paths},
  companion note, 2026.

\bibitem{Q_SCHUTZ}
  R.~Langer,
  \emph{The $Q$-Sch\"utzenberger Identity},
  companion note, 2026.

\bibitem{UNDERLYING_SET}
  R.~Langer,
  \emph{The Underlying Set: Cylindric Plane Partitions,
  ALCDs, and the Ground Set for the Involution},
  companion note, 2026.

\bibitem{langer-thesis}
  R.~Langer,
  \emph{Cylindric Partitions and Bijective Proofs},
  PhD thesis, University of Melbourne.
\end{thebibliography}

\end{document}
